\section{Superset Extensions Beyond Mermaid}
\label{sec:superset-extensions}

The compiler is not merely Mermaid-compatible---it is a \textbf{superset} that extends
Mermaid's capabilities through well-defined extension mechanisms. Extensions use Mermaid's
own frontmatter/directive hooks plus new top-level keywords, ensuring that the base Mermaid
grammar remains parseable by standard Mermaid renderers (with graceful degradation of
extended features).

\subsection{Extension Mechanisms}

\begin{description}
  \item[Frontmatter extensions] The YAML frontmatter block (\texttt{---}$\ldots$\texttt{---})
        carries our theme selection, animation directives, and composition metadata.
        Standard Mermaid renderers ignore unknown frontmatter fields.

  \item[Init directive extensions] The \texttt{\%\%\{init: \{...\}\}\%\%} block accepts
        grammar-specific configuration that maps to our theme extensions.

  \item[New top-level keywords] For capabilities with no Mermaid equivalent
        (e.g., \texttt{poster}, \texttt{compose}), we introduce new diagram-type keywords.
        These are clearly outside Mermaid's grammar and signal superset-only features.

  \item[Inline annotations] Extended comment syntax (\texttt{\%\% @directive}) for
        per-element hints (animation targets, cross-diagram references).
\end{description}

\subsection{Composition / Posters}

A new \texttt{poster} keyword enables multi-diagram compositions directly in the DSL:

\begin{lstlisting}[language=yaml,caption={Poster composition in superset DSL}]
---
theme: professional-dark
layout: grid 2x2
---
poster "RAG Architecture"
  cell [0,0]: flowchart LR
    A[Query] --> B[Retriever] --> C[Generator]
  cell [0,1]: sequence
    User->>API: request
    API->>DB: lookup
  cell [1,0]: mindmap
    root(Components)
      Retriever
      Generator
      Router
  cell [1,1]: stat
    value: 99.2%
    label: Accuracy
\end{lstlisting}

This maps to the Composition IR (Section~\ref{sec:composition}) and produces grid-based
multi-panel posters.

\subsection{Rich Theming}

Superset theming goes beyond Mermaid's color overrides:

\begin{lstlisting}[language=yaml,caption={Rich theme configuration via frontmatter}]
---
theme: corporate-blue
themeVariables:
  fontFamily: "Inter, system-ui"
  fontSize: 14
  borderRadius: 8
  nodePadding: 16
  lineWidth: 2
darkMode: true
---
\end{lstlisting}

Our theme system (Section~\ref{sec:themes}) provides full typographic, geometric, and
chromatic control---not just color token overrides.

\subsection{Animation}

Declarative animation hints attached to elements:

\begin{lstlisting}[language=yaml,caption={Animation extension}]
---
animate:
  edges: dashflow
  duration: 2s
---
flowchart LR
    A --> B --> C
\end{lstlisting}

Animation compiles to SVG SMIL (Section~\ref{sec:animation}) with raster backends
rendering static resting frames.

\subsection{Cross-Diagram References}

Elements in one diagram can reference elements in another within the same poster:

\begin{lstlisting}[language=yaml,caption={Cross-diagram reference annotation}]
---
refs:
  - from: flow.nodeC
    to: sequence.participant.API
    style: dashed-arrow
---
\end{lstlisting}

References resolve at the composition layer and render as overlay connectors.

\subsection{Icons and Images in Nodes}

Nodes can embed icons from a registry or external images:

\begin{lstlisting}[language=yaml,caption={Icon extension in flowchart}]
flowchart LR
    A["{icon:aws-lambda}" Lambda] --> B["{icon:database}" Store]
\end{lstlisting}

The icon registry is part of the kernel; themes control icon styling (color, size).

\subsection{The Structured IR as Agent API}

The superset's most significant extension is invisible in the DSL: the structured IR path.
Agents bypass the text parser entirely and emit Domain IR directly as JSON/YAML.
This is not an ``export format''---it is a first-class input API
(Section~\ref{sec:frontend}).

\subsection{Graceful Degradation}

When a superset file is rendered by a standard Mermaid renderer:
\begin{itemize}
  \item Frontmatter extensions are silently ignored (Mermaid skips unknown fields).
  \item New keywords (\texttt{poster}, etc.) produce a parse error---clearly signaling
        that the file requires our renderer.
  \item Inline annotations in comments (\texttt{\%\% @...}) are ignored as comments.
\end{itemize}

This ensures that the base diagram content remains portable; only superset features
require our compiler.
